% =====================================================
%  PLANTILLA DE INFORME
%  Compilar con: pdflatex (2 veces, para índice y refs)
% =====================================================
\documentclass[12pt,a4paper]{article}

	\usepackage{xcolor}
	\usepackage{listings}
	\usepackage{tcolorbox}
	\definecolor{main}{RGB}{145,145,145}
	\usepackage{lscape}
	
	% ---------- Codificación e idioma ----------
	\usepackage[utf8]{inputenc}
	\usepackage[T1]{fontenc}
	\usepackage{lmodern}
	\usepackage[spanish,es-tabla]{babel}   % "Tabla" en vez de "Cuadro"
	\decimalpoint                          % punto decimal (Perú/LatAm)
	
	% ---------- Formato de página ----------
	\usepackage[margin=2cm]{geometry} % 2.5cm o 2cm
	\usepackage{setspace}
	\onehalfspacing
	\usepackage{microtype}                 % mejor justificación del texto
	\setlength{\parindent}{1.25cm}
	\setlength{\parskip}{4pt}
	
	% ---------- Encabezado y pie de página ----------
	\usepackage{fancyhdr}
	\pagestyle{fancy}
	\fancyhf{}
	\fancyhead[L]{\small\itshape \tituloinforme}
	%\fancyhead[R]{\small\leftmark}
	\fancyhead[R]{\small Monitoreo de participación semanal}
	\fancyfoot[C]{\thepage}
	\renewcommand{\headrulewidth}{0.4pt}
	
	% ---------- Títulos de sección ----------
	\usepackage{titlesec}
	\titleformat{\section}{\Large\bfseries}{\thesection.}{0.6em}{}
	\titleformat{\subsection}{\large\bfseries}{\thesubsection.}{0.6em}{}
	
	% ---------- Matemáticas, figuras, tablas, código ----------
	\usepackage{amsmath,amssymb}
	\usepackage{graphicx}
	\usepackage{booktabs}
	\usepackage{caption}
	\captionsetup{font=small,labelfont=bf}
	\usepackage{listings}
	\lstset{basicstyle=\ttfamily\small,breaklines=true,frame=single}
	
	% ---------- Enlaces y PDF ----------
	\usepackage[hidelinks]{hyperref}
	
	% ---------- Datos del informe (editar aquí) ----------
	\newcommand{\tituloinforme}{Caso Práctico}
	\newcommand{\subtituloinforme}{Monitoreo de participación semanal}
	\newcommand{\autorinforme}{Jhon R. Ordoñez}
	\newcommand{\institucion}{Aprendo +}
	\newcommand{\fechainforme}{\today}
	
	\hypersetup{pdftitle={\tituloinforme}, pdfauthor={\autorinforme}}

	% ---------- Soporte de tildes y ñ dentro del código (con utf8) ----------
	\lstset{
		literate=
		{á}{{\'a}}1 {é}{{\'e}}1 {í}{{\'i}}1 {ó}{{\'o}}1 {ú}{{\'u}}1
		{Á}{{\'A}}1 {É}{{\'E}}1 {Í}{{\'I}}1 {Ó}{{\'O}}1 {Ú}{{\'U}}1
		{ñ}{{\~n}}1 {Ñ}{{\~N}}1
	}
	\lstdefinestyle{tsql}{
		language=SQL,
		morekeywords={use, go, if, begin, end, if, top},
		basicstyle=\ttfamily\small,
		keywordstyle=\color{blue}\bfseries,
		commentstyle=\color{green!50!black}\itshape,
		stringstyle=\color{red!70!black},
		numbers=none, % none or left 
		numberstyle=\tiny\color{gray},
		frame = none,
		breaklines=true,
		showstringspaces=false,
		captionpos=b,
		tabsize=4,
		morecomment=[l]{--},
		morecomment=[s]{/*}{*/},
		morestring=[b]'
	}

% =====================================================
\begin{document}
	
	% ---------- Portada ----------
	\begin{titlepage}
		\centering
		\vspace*{2cm}
		{\huge\scshape \institucion\par}
		\vspace{3cm}
		\rule{\textwidth}{1.2pt}\par\vspace{0.6cm}
		{\Huge\bfseries \tituloinforme\par}
		\vspace{0.4cm}
		{\Large \subtituloinforme\par}
		\vspace{0.6cm}
		\rule{\textwidth}{1.2pt}\par
		\vspace{0.6cm}
			\begin{abstract}
				\noindent La participación de las familias se mantiene en torno al 33\% en la semana 6, pero ese
				promedio esconde un problema concentrado en tres instituciones rurales de UGEL C (IE-C05,
				IE-C06 e IE-C07), donde la participación cayó de aproximadamente 37\% a 7\%.
			 	Estas instituciones
				incluyen los cinco grupos con mayor caída y 40 de las 142 familias que dejaron de participar,
				por lo que es ahí donde conviene actuar primero. Además, la baja de UGEL B en la semana 3 (21.8\%)
				no refleja un desinterés de las familias, sino la ausencia de registros los días 2 y 3 de
				septiembre, por lo que esa cifra no es comparable y debe leerse como dato incompleto.
			\end{abstract}
		\vfill
		{\Large \autorinforme\par}
		\vspace{0.5cm}
		{\large \fechainforme\par}
		\vspace{1cm}
	\end{titlepage}
	
	% ---------- Resumen ----------
	%\begin{abstract}
	%	Escriba aquí un resumen breve del informe: objetivo, método
	%	y principales resultados ...
	%\end{abstract}
	
	%\newpage
	%\tableofcontents
	%\newpage
	
	% ---------- Section 1 ----------
	\section{SQL}
		

		% ---------- Subsection 1.1 ----------
		\subsection{Tasa de participación semanal}
			
			% Entregable 1.1
			\noindent \textbf{Pregunta 1.1:} La tasa de participación semanal por UGEL y en total, para las 6 semanas.
			
			
			\noindent Para poder entender la siguiente tabla, recordemos la definición de la métrica principal.
		
			\begin{center}
				\setlength{\fboxsep}{6pt}
				\setlength{\fboxrule}{0.8pt}
				\fbox{%
					\parbox{0.80\textwidth}{%
					\centering
					\textbf{Tasa de participación semanal} = \textcolor{blue}{familias distintas con al menos una 
					respuesta con evidencia (imagen, foto, audio o video) en las actividades de la semana} \textbf{÷} \textcolor{teal}{familias 
					registradas hasta el fin de esa semana}.
				}}
			\end{center}
			
			\noindent Una persona cuenta una sola vez aunque tenga más de un registro. 
				
			% Tabla 1
			\begin{table}[h]
				\centering
				\caption{Tasa de participación semanal, según ugel (\%)}
				\label{tab: tasa}
				\begin{tabular}{ccccc}
					\toprule
					Semana & UGEL A & UGEL B & UGEL C & Total \\
					\midrule
					1 & 38.4 & 37.2 & 31.0 & 35.5 \\
					2 & 43.3 & 36.0 & 32.3 & 37.1 \\
					3 & 42.4 & 21.8 & 36.0 & 33.4 \\
					4 & 48.1 & 41.0 & 36.6 & 41.8 \\
					5 & 42.1 & 39.5 & 28.8 & 36.8 \\
					6 & 40.9 & 34.2 & 24.2 & 33.0 \\
					\bottomrule
				\end{tabular}
			\end{table}
			
		\noindent La participación total se mantuvo entre 33.0 \% y 41.8 \% durante las seis semanas,
		con un pico en la semana 4 (41.8 \%, el mayor valor en las tres UGEL) y una caída posterior
		hasta 33.0 \% en la semana 6. UGEL A es la de mejor desempeño y mayor estabilidad (38.4 \%-48.1\%),
		UGEL B es la más volátil, con un descenso marcado en la semana 3 (21.8 \%, el mínimo de toda la tabla)
		seguido de recuperación y nueva baja gradual, y UGEL C tiene el rendimiento más bajo y un deterioro
		reciente, pues pasó de 36.6 \% en la semana 4 a 24.2 \% en la semana 6. La brecha entre UGEL A y C se
		amplió de 7.4 a 16.7 puntos porcentuales, por lo que se recomienda identificar qué favoreció la
		participación en la semana 4, investigar la caída de UGEL B en la semana 3 y priorizar el acompañamiento 
		en UGEL C, tomando las prácticas de UGEL A como referencia.
			
		% ---------- Subsection 1.2 ----------
		\subsection{Cinco grupos con mayor caida}

			% Entregable 1.2
			\noindent \textbf{Pregunta 1.2:} Los 5 grupos con mayor caída de participación entre el promedio de las semanas 1 a 
			4 y el promedio de las semanas 5 y 6. 
		
			% Tabla 2
			\begin{table}[h]
				\centering
				\caption{Los cinco grupos con mayor caída}
				\label{tab: tasa}
				\begin{tabular}{ccccccc}
					\toprule
					N° & Grupo & UGEL  & Prom. S1-S4(\%) & Prom. S5-S6(\%) & Caída (pp) & Familias \\
					\midrule
					1 & G-C06-5 & UGEL C & 44.2 & 9.6  & 34.6 & 26 \\
					2 & G-C07-5 & UGEL C & 38.3 & 5.9  & 32.4 & 17 \\
					3 & G-C07-4 & UGEL C & 37.1 & 8.8  & 28.3 & 17 \\
					4 & G-C06-4 & UGEL C & 35.5 & 8.3  & 27.1 & 18 \\
					5 & G-C05-5 & UGEL C & 38.3 & 12.5 & 25.8 & 24 \\
					\bottomrule
				\end{tabular}
			\end{table}
		
		\noindent Cinco grupos de UGEL C, todos de tres instituciones rurales (IE-C05, IE-C06 e IE-C07), pasaron de 
		una participación de 35\% a 44\% en las semanas 1 a 4 a solo 6\% a 13\% en las semanas 5 y 6, una pérdida de 
		26 a 35 puntos porcentuales. La mayor caída es G-C06-5, que bajó de 44.2\% a 9.6\%. Estos grupos suman 102 familias,
		cerca del 10\% de las 1,023 registradas, y por eso la tasa total solo bajó de 36.8\% a 33.0\%, ya que el 
		promedio oculta el problema. En cuatro de los cinco grupos las docentes no escribieron nada en la semana 5,
		y en el quinto casi nada, lo que apunta a un problema operativo más que a desinterés de las familias, aunque 
		con estos datos no podemos confirmar la causa. Los grupos son pequeños (17 a 26 familias), pero la caída se
		repite en los cinco, así que no parece azar. Esta semana conviene pedir a Rocket Learning el estado de
		esos grupos, llamar a las docentes de las tres instituciones y definir un plan de recuperación.
		
		
		% ---------- Subsection 1.3 ----------
		\subsection{Familias que dejan de participar}
			
			% Entregable 1.3
			\noindent \textbf{Pregunta 1.3:} Cuántas familias participaron con evidencia en al menos 2 semanas distintas, de las 
			semanas 1 a 4 y no participaron con evidencia en las semanas 5 y 6, y en qué grupos se 
			concentran. 
			
			% Tabla 3
			\begin{table}[h]
				\centering
				\caption{Familias que dejan de participar}
				\label{tab: tasa}
				\begin{tabular}{cccccc}
					\toprule
					Grupo & UGEL & Familias que dejan & Familias del grupo & \% del grupo & \% acum. del total  \\
					\midrule
					G-C05-5 & C & 10 & 24 & 41.7 &  7.0 \\
					G-C06-5 & C & 10 & 26 & 38.5 & 14.1 \\
					G-A04-5 & A & 5  & 21 & 23.8 & 17.6 \\
					G-B02-5 & B & 5  & 17 & 29.4 & 21.1 \\
					G-C05-4 & C & 5  & 25 & 20.0 & 24.6 \\
					G-C06-4 & C & 5  & 18 & 27.8 & 28.2 \\
					G-C07-4 & C & 5  & 17 & 29.4 & 31.7 \\
					G-C07-5 & C & 5  & 17 & 29.4 & 35.2 \\
					\vdots & \vdots & \vdots  & \vdots  & \vdots  & \vdots  \\
					G-C04-5& C &  1  & 19 & 5.3 & 100.0 \\
					\bottomrule
				\end{tabular}
			\end{table}
			
		\noindent Entre las 142 familias que dejaron de participar, la pérdida se concentra más en unos pocos grupos
		que en el conjunto. Son familias que participaban con evidencia en al menos dos de las primeras
		cuatro semanas y no lo hicieron en las semanas 5 y 6. Dos grupos, G-C05-5 y G-C06-5, perdieron cada
		uno 10 familias, cerca de 4 de cada 10 de las suyas (41.7\% y 38.5\%). Juntos explican el 14.1\% del 
		total, y los ocho grupos más afectados suman el 35.2\%. Seis de esos ocho grupos están en UGEL C.
		Otros 39 grupos pierden entre 1 y 4 familias cada uno, lo que parece un desgaste disperso, sin una 
		causa común visible. Por eso conviene atender primero los grupos de UGEL C, donde la pérdida es
		alta en términos relativos (entre 20\% y 42\% del grupo), antes que aplicar medidas generales.
		
		% ---------- Subsection 1.4 ----------
		\subsection{Controles de calidad}
		
			% Entregable 1.4
			\noindent \textbf{Pregunta 1.4:} Consultas de control de calidad que cuantifiquen los problemas que encuentres en 
			los datos. 
			
			\begin{enumerate}
				\item La variable tipo tiene variantes en sus categorías Foto, 
				      FOTO (con espacio al final), imagen, Audio, Video y Reaccion; estas variantes
				      afectan a 391 filas (representa el 5.5 \% de 7077 filas), de las cuales 351 filas son 
				      categorias consideradas como evidencia y 40 no se consideran evidencia.
				      La consulta lo puede ver en el \textcolor{blue}{anexo \ref{anexo 5}}            
				\item En la tabla participantes una misma persona (telefono\_hash) con 2 participante\_id, siempre
					  en el mismo grupo, estos duplicados afecta a 60 filas = 30 personas (30 sobrantes).
					  La consulta lo puede ver en el \textcolor{blue}{anexo \ref{anexo 5}}
				\item En la tabla interacciones se tiene duplicado exacto, donde mismo participante, actividad, instante y tipo con distinto
				      interaccion\_id. Esos duplicados generan 137 filas sobrantes que representa el 1.9\% de 7077.
				      La consulta lo puede ver en el \textcolor{blue}{anexo \ref{anexo 5}}.
				\item Respuestas fechada antes del envío de la actividad (hora de
				      Lima; hasta 34.6 h antes), donde 50 filas son afectadas que representa el 0.7\% de 7077, 
				      de las cuales 31 filas con evidencia. 
				      La consulta lo puede ver en el \textcolor{blue}{anexo \ref{anexo 5}}.
				\item Se identifica un vacío de datos en la semana 3 para la UGEL B, específicamente en las actividades 13
				      y 14, correspondientes al 2 y 3 de septiembre, respectivamente. En ambas actividades no se registraron
				      interacciones durante las fechas señaladas, por lo que el número de interacciones es igual a cero.
					  La consulta lo puede ver en el \textcolor{blue}{anexo \ref{anexo 5}}.
			\end{enumerate}

		
		
	% ---------- Section 2 ----------
	\newpage
	\section{Calidad de datos}
		
		% Entregable 1.4
		\noindent \textbf{Pregunta 2:} Lista priorizada de los problemas que 
		encontraste, su impacto en la métrica, qué corregirías tú y qué le pedirías al equipo de 
		Rocket Learning. 
		
			\begin{enumerate}
				\item Durante la semana 3 (2-3 de septiembre) se detectó un hueco en los datos de la UGEL B, 
				      con 0 registros en ACT013-014, esto distorsiona la métrica: la UGEL B marca 21.8\% en la
				      semana 3 frente a 36.0\% y 41.0\% en semana 2 y semana 4, y el total de semana 3 queda en 33.4\%,
				      lo que sesga el promedio de la semana 1 a 4 de los 16 grupos. Mi propuesta es no imputar los valores
				      faltantes, sino marcar la semana 3 del UGEL B como “dato incompleto”. A Rocket Learning le pediría que identifique la causa
				      (plataforma, ingesta o WhatsApp), que vuelva a revisar los datos faltantes y que 
				      comparta los hallazgos.
				\item Se detectó un problema de estandarización en el campo “tipo”: 391 filas no están normalizadas,
				      de las cuales 351 contienen evidencia (el 10.6\% del total de evidencias). Dejar estos valores
				      sin normalizar afecta la métrica. Para corregirlo, aplicaría
				      lower(ltrim(rtrim())) para limpiar el texto, unificaría “imagen” como “foto” y crearía una tabla
				      de dominio con los valores válidos. A Rocket Learning le pediría una lista cerrada de tipos,
				      validada desde el origen, para evitar que el problema se repita.
				\item Se detectaron 30 personas registradas con dos IDs distintos (60 filas en total), que comparten 
				      el mismo teléfono y grupo. Esto infla el denominador de la métrica. Para corregirlo, definiría
				      a la persona por su telefono\_hash y conservaría como alta
				      válida la más antigua. A Rocket Learning le pediría un ID de persona estable, que explique por
				      qué las personas se re-registran y que aplique la deduplicación desde el origen.
				\item Se detectaron 50 respuestas registradas antes de la hora 
				      de envío, según la hora de Lima, con diferencias de hasta 34.6 horas.
				      A Rocket Learning le pediría que revise el reloj o huso horario y la asignación del actividad\_id.
				\item Se detectaron 137 filas duplicadas exactas. Esto no afecta la tasa, porque esta
				      cuenta personas únicas, pero sí infla los conteos de mensajes.
			\end{enumerate}
			
			
	
	% ---------- Section 3 ----------
	\newpage 
	\section{Alerta}
		
		% Entregable 1.4
		\noindent \textbf{Pregunta 3:} Un mensaje al CEO y a la PMO Lead, como si fuera el 
		lunes 28, sobre lo que muestra el reporte. Debe poder leerse en un minuto por alguien no 
		técnico y decir qué sabes, qué no sabes y qué harías esta semana. \\ [3mm]
	
		\noindent \textbf{Asunto: Participación semana 6 (lunes 28)}
		
		\noindent Estimado (a),
		
		\noindent \textbf{Lo que sabemos:} el 33\% de las familias envió imagen, foto, audio o
		video esta semana, mientras que la semana anterior fue 36.8\%. La caída
		se concentra en las instituciones educativas rurales 5, 6 y 7
		del UGEL C de aproximadamente 37\% a 7\%, y allí casi no
		escriben las docentes.
		
		\noindent \textbf{Lo que no sabemos:} si es desinterés, un problema técnico del grupo de WhatsApp o un error de datos. Además, UGEL
		B no tiene registros el 2 y 3 de septiembre; su semana 3 no es comparable.
		
		\noindent \textbf{Lo que haremos esta semana:} hoy pedimos a Rocket Learning confirmar el hueco y el estado de esos grupos;
		mañana llamamos a las docentes; el viernes informo causa y plan.
		
		\noindent Atentamente,
		
		\noindent Jhon R. Ordoñez
		
	% ---------- Section 4 ----------
	\newpage
	\section{Automatización}
		
		% Entregable 1.4
		\noindent \textbf{Pregunta 4:} Cómo dejarías este reporte 
		corriendo todos los días y cada semana sobre AWS: qué se ejecuta y cuándo, qué 
		validaciones incluye, qué pasa si algo falla y quién se entera. Si no conoces AWS, propón el 
		servicio equivalente. También: Indica tus supuestos sobre cómo llegan los datos.
	
		% Diagrama básico
		\begin{center}
			\setlength{\fboxsep}{6pt}
			\setlength{\fboxrule}{0.8pt}
			
			\fbox{\textcolor{darkgray}{\textbf{1. Obtener data}}}
			
			Recibir los archivos .csv en una carpeta con fecha (Manual al inicio)
			
			\vspace{2mm}
			$\downarrow$
			\vspace{2mm}
			
			\fbox{\textcolor{blue}{\textbf{2. Cargar a SQL}}}
			
			Cargar los archivos tal cual a tablas crudas (bulk insert)
			
			\vspace{2mm}
			$\downarrow$
			\vspace{2mm}
			
			\fbox{\textcolor{blue}{\textbf{3. Validar}}}
			
			Correr los controles de calidad y guardar el resultado
			
			\fbox{%
				\parbox{0.80\textwidth}{%
					\centering
					\textcolor{red}{
						Si un control de calidad falla, se suspende la publicación y se comunica
						de inmediato la incidencia al líder del proyecto. Mientras se espera una
						respuesta, se investiga la causa y se evalúan posibles soluciones.
						La publicación se reanuda únicamente después de corregir y validar
						satisfactoriamente la incidencia.  
					}
			}}
			
			\vspace{2mm}
			$\downarrow$
			\vspace{2mm}
			
			\fbox{\textcolor{blue}{\textbf{4. Tablas de análisis}}}
			
			Calcular la métrica una sola vez (Tablas auxiliares y consultas)
			
			\vspace{2mm}
			$\downarrow$
			\vspace{2mm}
			
			\fbox{\textcolor{teal}{\textbf{5. Power BI}}}
			
			Leer las tablas ya listas (Conexión a SQL Server)
			
			\vspace{2mm}
			$\downarrow$
			\vspace{2mm}
			
			\fbox{\textcolor{teal}{\textbf{6. Dashboard} (3 páginas)}}
			
			Mostrar el resultado (Participación, alertas y calidad de datos)
			
		\end{center}
	
	
	% ---------- Anexos ----------
	\newpage	
	\begin{center}
		\vspace*{10cm}
		{\Huge \textbf{ANEXO}}
	\end{center}
	\newpage 
	
	\appendix
	% Appendix 1	
	%\begin{landscape}
	%\end{landscape}
	\section{Consultas SQL para generar tablas auxiliares}\label{anexo 1}
	
	    \noindent Para el desarrollo y ejecución de las consultas se utilizó Microsoft SQL Server como motor
		de base de datos, Transact-SQL (T-SQL) como lenguaje de consulta y SQL Server Management Studio (SSMS)
		como entorno de administración y ejecución.
		
		% Tabla 1
		\begin{tcolorbox}[colback= main!40, colframe=main!40]
			\lstinputlisting[style=tsql, basicstyle=\fontsize{7}{5}\selectfont\ttfamily]{queries/table_1.sql}
		\end{tcolorbox}
		% Tabla 2
		\begin{tcolorbox}[colback= main!40, colframe=main!40]
			\lstinputlisting[style=tsql, basicstyle=\fontsize{7}{5}\selectfont\ttfamily]{queries/table_2.sql}
		\end{tcolorbox}
		% Tabla 3
		\begin{tcolorbox}[colback= main!40, colframe=main!40]
			\lstinputlisting[style=tsql, basicstyle=\fontsize{7}{5}\selectfont\ttfamily]{queries/table_3.sql}
		\end{tcolorbox}
		% Tabla 4
		\begin{tcolorbox}[colback= main!40, colframe=main!40]
			\lstinputlisting[style=tsql, basicstyle=\fontsize{7}{5}\selectfont\ttfamily]{queries/table_4.sql}
		\end{tcolorbox}
	
	% Appendix 2
	\section{Consulta SQL de la parte 1.1}\label{anexo 2}
		
		\noindent Para el desarrollo y ejecución de las consultas se utilizó Microsoft SQL Server como motor
		de base de datos, Transact-SQL (T-SQL) como lenguaje de consulta y SQL Server Management Studio (SSMS)
		como entorno de administración y ejecución.
		
		% Consulta 1
		\begin{tcolorbox}[colback = main!40, colframe = main!40]
			\lstinputlisting[style=tsql, basicstyle=\fontsize{4.4}{4}\selectfont\ttfamily]{queries/parte1.1_query_1.sql}
		\end{tcolorbox}
		
	% Appendix 3
	\section{Consulta SQL de la parte 1.2}\label{anexo 3}
		
		\noindent Para el desarrollo y ejecución de las consultas se utilizó Microsoft SQL Server como motor
		de base de datos, Transact-SQL (T-SQL) como lenguaje de consulta y SQL Server Management Studio (SSMS)
		como entorno de administración y ejecución.
		
		% Consulta 2
		\begin{tcolorbox}[colback = main!40, colframe = main!40]
			\lstinputlisting[style=tsql, basicstyle=\fontsize{4.4}{4}\selectfont\ttfamily]{queries/parte1.2_query_1.sql}
		\end{tcolorbox}
		
	% Appendix 4
	\section{Consulta SQL de la parte 1.3}\label{anexo 4}
		
		\noindent Para el desarrollo y ejecución de las consultas se utilizó Microsoft SQL Server como motor
		de base de datos, Transact-SQL (T-SQL) como lenguaje de consulta y SQL Server Management Studio (SSMS)
		como entorno de administración y ejecución.
		
		% Consulta 3
		\begin{tcolorbox}[colback = main!40, colframe = main!40]
			\lstinputlisting[style=tsql, basicstyle=\fontsize{4.4}{4}\selectfont\ttfamily]{queries/parte1.3_query_1.sql}
		\end{tcolorbox}
	
	% Appendix 5
	\section{Consulta SQL de la parte 1.4}\label{anexo 5}
		
		\noindent Para el desarrollo y ejecución de las consultas se utilizó Microsoft SQL Server como motor
		de base de datos, Transact-SQL (T-SQL) como lenguaje de consulta y SQL Server Management Studio (SSMS)
		como entorno de administración y ejecución.
		
		% Consulta 4
		\begin{tcolorbox}[colback = main!40, colframe = main!40]
			\lstinputlisting[style=tsql, basicstyle=\fontsize{4.4}{4}\selectfont\ttfamily]{queries/parte1.4_query_1.sql}
		\end{tcolorbox}
		% Consulta 5
		\begin{tcolorbox}[colback = main!40, colframe = main!40]
			\lstinputlisting[style=tsql, basicstyle=\fontsize{4.4}{4}\selectfont\ttfamily]{queries/parte1.4_query_2.sql}
		\end{tcolorbox}
		% Consulta 6
		\begin{tcolorbox}[colback = main!40, colframe = main!40]
			\lstinputlisting[style=tsql, basicstyle=\fontsize{4.4}{4}\selectfont\ttfamily]{queries/parte1.4_query_3.sql}
		\end{tcolorbox}
		% Consulta 7
		\begin{tcolorbox}[colback = main!40, colframe = main!40]
			\lstinputlisting[style=tsql, basicstyle=\fontsize{4.4}{4}\selectfont\ttfamily]{queries/parte1.4_query_4.sql}
		\end{tcolorbox}
		% Consulta 8
		\begin{tcolorbox}[colback = main!40, colframe = main!40]
			\lstinputlisting[style=tsql, basicstyle=\fontsize{4.4}{4}\selectfont\ttfamily]{queries/parte1.4_query_5.sql}
		\end{tcolorbox}

	\label{key}
	% Appendix 5
	\section{Definiciones y supuestos usados}
	\noindent Participación es igual a familia con mayor o igual a una interacción de tipo
	foto, imagen, audio o video en alguna actividad de la semana;
	solo rol familia; persona es igual a telefono\_hash (una persona cuenta una vez aunque tenga 2 identificadores); 
	denominador es igual a familias con fecha\_alta menor o igual a domingo de esa
	semana; la semana es la de la actividad respondida; texto, reacción y reenvío no son evidencia;
	identificadores no registrados se excluyen. El analista de tecnología de Aprendo+ confirmó que  
	imagen es igual a foto; cierre semanal es igual a domingo, con todas las familias registradas hasta ese día;
	ninguna familia sale de los grupos (el denominador nunca baja). Zona horaria de Perú (UTC-5, sin
	horario de verano); las actividades salen a las 18:00
	Lima (23:00 UTC) y cada respuesta va a la semana de su actividad.
	
	
\end{document}







